\section{正弦と余弦を交換する}\label{節:正弦余弦の交換}

区間 $\left[0,\frac{\pi}{2}\right]$ での反転は，$x$ を $\frac{\pi}{2}-x$ に変える操作である。このとき
\[
\sin\left(\frac{\pi}{2}-x\right)=\cos x,\qquad
\cos\left(\frac{\pi}{2}-x\right)=\sin x
\]
なので，正弦と余弦が入れ替わる。分子と分母に両者が現れる場合，入れ替えた式との和を調べるとよい。

\begin{example}{分母を変えずに分子を交換する}
次の定積分を求めよ。
\[
I=\int_0^{\frac{\pi}{2}}
\frac{\sin^4x}{\sin^4x+\cos^4x}\dd x.
\]
\end{example}

$\sin x$ と $\cos x$ は同時に零にならないので，分母は区間全体で正である。反転すると分子は $\cos^4x$ に変わり，分母は変わらない。したがって
\[
I=\int_0^{\frac{\pi}{2}}
\frac{\cos^4x}{\sin^4x+\cos^4x}\dd x.
\]
二つを加えれば，分子の和が分母そのものになる。
\[
2I=\int_0^{\frac{\pi}{2}}1\dd x=\frac{\pi}{2},
\qquad I=\frac{\pi}{4}.
\]
$\sin^4x$ の次数を下げたり，分母を展開したりする前に，分子の足し合わせで分母ができることに気づくと計算を減らせる。

この関係は四乗に特有ではない。$q$ を $[0,1]$ 上の連続関数とし，$q(\sin x)+q(\cos x)$ が区間内で零にならなければ，
\[
\int_0^{\frac{\pi}{2}}
\frac{q(\sin x)}{q(\sin x)+q(\cos x)}\dd x
 =\frac{\pi}{4}
\]
である。たとえば $q(t)=t^2+1$ でも同じ値になる。重要なのは指数の値ではなく，反転で分子の二つの候補が交換されることである。

図\ref{図:相補重み}は，例題の被積分関数を $w(x)$ としたグラフである。左右を反転したグラフ $w(\frac{\pi}{2}-x)$ と，同じ $x$ における高さを足すと一になる。それぞれの積分が等しいので，全体の長方形の面積を二等分している。

\begin{figure}[htbp]
\centering
\begin{tikzpicture}[x=3.5cm,y=2.3cm,>=Stealth]
\fill[blue!10](0,0)--plot[domain=0:1.570796,samples=100](\x,{pow(sin(\x r),4)/(pow(sin(\x r),4)+pow(cos(\x r),4))})--(1.570796,0)--cycle;
\draw[->](-0.08,0)--(1.73,0)node[right]{$x$};
\draw[->](0,-0.08)--(0,1.2)node[above]{$y$};
\draw[dashed](0,1)--(1.570796,1)--(1.570796,0);
\draw[dashed](0.785398,0)--(0.785398,1);
\draw[blue!75!black,thick]plot[domain=0:1.570796,samples=100](\x,{pow(sin(\x r),4)/(pow(sin(\x r),4)+pow(cos(\x r),4))});
\draw[orange!80!black,thick,dashed]plot[domain=0:1.570796,samples=100](\x,{pow(cos(\x r),4)/(pow(sin(\x r),4)+pow(cos(\x r),4))});
\node[left]at(0,1){$1$};\node[below]at(0,0){$0$};
\node[below]at(0.785398,0){$\frac{\pi}{4}$};
\node[below]at(1.570796,0){$\frac{\pi}{2}$};
\node[blue!75!black]at(1.15,0.65){$w(x)$};
\node[orange!80!black]at(0.48,1.26){$w(\frac{\pi}{2}-x)$};
\end{tikzpicture}
\caption{反転した二つの重みは，同じ場所で足すと一になる。}\label{図:相補重み}
\end{figure}

